\documentclass[11pt,a4paper]{article}
\usepackage[T1]{fontenc}
\usepackage[utf8]{inputenc}
\usepackage{lmodern}
\usepackage{amsmath,amssymb,amsthm,mathtools,mathrsfs}
\usepackage{graphicx,float,xcolor,multicol}
\usepackage[shortlabels]{enumitem}
\usepackage[margin=27mm]{geometry}
\usepackage{microtype,needspace,placeins}
\usepackage{tikz,tikz-cd}
\usetikzlibrary{arrows.meta,calc,positioning,patterns,decorations.pathreplacing}
\usepackage[colorlinks=true,linkcolor=teal,urlcolor=teal,citecolor=teal]{hyperref}
\setlength{\parindent}{0pt}
\setlength{\parskip}{.6em}
\setcounter{secnumdepth}{0}
\allowdisplaybreaks
\raggedbottom
\providecommand{\cross}{\times}
\DeclareUnicodeCharacter{2020}{\ensuremath{\dagger}}
\DeclareMathOperator{\supp}{supp}
\DeclareMathOperator*{\esssup}{ess\,sup}
\providecommand{\chapter}{\section}
\newenvironment{tool}[2]{\subsection*{Tool #1 --- #2}}{}
\newcommand{\ToolRef}[1]{Tool #1}
\newcommand{\FigureTag}[2]{\textit{Figure #1}}
\newcommand{\Result}[1]{\subsection*{#1}}
\newcommand{\Application}[1]{\subsubsection*{Application\if\relax\detokenize{#1}\relax\else\space--- #1\fi}}
\newcommand{\ProofHeading}{\par\textit{Proof.}\space}
\newcommand{\ProofOf}[1]{\par\textit{Proof of #1.}\space}
\newcommand{\RemarkHeading}{\par\textbf{Remark.}\space}
\newcommand{\NoteHeading}{\par\textbf{Note.}\space}
\newcommand{\ExerciseHeading}{\par\textbf{Exercise.}\space}

% Rendering support only; mathematical text is in each independent post.
\usepackage{booktabs,array,longtable}
\definecolor{HandoutBlue}{HTML}{2F6690}
\definecolor{HandoutNavy}{HTML}{17365D}
\newtheorem{theorem}{Theorem}
\newtheorem{lemma}[theorem]{Lemma}
\newtheorem{proposition}[theorem]{Proposition}
\newtheorem{corollary}[theorem]{Corollary}
\newtheorem{conjecture}[theorem]{Conjecture}
\newtheorem{definition}{Definition}
\newtheorem{example}{Example}
\newtheorem{namedtoolcounter}{Tool}
\newenvironment{namedtool}[1]{\begin{namedtoolcounter}[#1]}{\end{namedtoolcounter}}
\newtheorem*{claim}{Claim}
\newtheorem*{remark}{Remark}
\newtheorem*{exercise}{Exercise}
\newtheorem*{question}{Question}
\newtheorem*{observation}{Observation}
\newcommand{\SourceChapter}[2]{\section*{#2}}
\newcommand{\Lecture}[2]{\section*{Lecture #1: #2}}
\newcommand{\unclear}[1]{\textit{[unclear: #1]}}
\newcommand{\sourcegap}{\textit{[unfinished in the source]}}
\DeclareMathOperator{\dist}{dist}
\DeclareMathOperator{\id}{id}
\DeclareMathOperator{\Hol}{Hol}
\DeclareMathOperator{\Per}{Per}
\DeclareMathOperator{\Fix}{Fix}
\DeclareMathOperator{\diam}{diam}
\DeclareMathOperator{\Var}{Var}
\DeclareMathOperator{\Leb}{Leb}
\DeclareMathOperator{\Hom}{Hom}
\DeclareMathOperator{\End}{End}
\DeclareMathOperator{\Ann}{Ann}
\DeclareMathOperator{\Spec}{Spec}
\DeclareMathOperator{\Max}{Max}
\DeclareMathOperator{\Ob}{Ob}
\DeclareMathOperator{\Tor}{Tor}
\DeclareMathOperator{\Ass}{Ass}
\DeclareMathOperator{\Mor}{Mor}
\DeclareMathOperator{\Fun}{Fun}
\DeclareMathOperator{\Nat}{Nat}
\DeclareMathOperator{\Cone}{Cone}
\DeclareMathOperator{\MCG}{MCG}
\DeclareMathOperator{\Aut}{Aut}
\DeclareMathOperator{\Deck}{Deck}
\DeclareMathOperator{\SL}{SL}
\DeclareMathOperator{\PSL}{PSL}
\DeclareMathOperator{\Area}{Area}
\DeclareMathOperator{\im}{im}
\DeclareMathOperator{\PSh}{PSh}
\newcommand{\CRing}{\mathsf{CRings}}
\newcommand{\AMod}{A\text{-}\mathsf{Mod}}
\newcommand{\Ab}{\mathsf{Ab}}
\newcommand{\Z}{\mathbb Z}
\newcommand{\Q}{\mathbb Q}
\newcommand{\R}{\mathbb R}
\newcommand{\C}{\mathbb C}
\newcommand{\N}{\mathbb N}
\newcommand{\kk}{\mathbb k}
\renewcommand{\aa}{\mathfrak a}
\newcommand{\bb}{\mathfrak b}
\newcommand{\pp}{\mathfrak p}
\newcommand{\qq}{\mathfrak q}
\newcommand{\mm}{\mathfrak m}
\newcommand{\nn}{\mathfrak n}
\newcommand{\rad}{\mathrm r}
\newcommand{\cat}[1]{\mathsf{#1}}
\setcounter{secnumdepth}{0}
\allowdisplaybreaks

\graphicspath{{figures/lemma-book/}}
\begin{document}
\section*{Polynomial Roots, Identities, and Algebraic Systems}
% BLOG-CONTENT-BEGIN
\setcounter{theorem}{13}\begin{lemma}[Detecting nonreal roots]
If one is asked to prove that a polynomial has complex roots, focus on showing that the sum of the squares of its roots is negative.
\end{lemma}
\paragraph{Application.}
Show that $ax^4+bx^3+x^2+x+1=0$, where $a,b\in\mathbb R$ and $a\ne0$, has complex roots.
Let $\alpha_1,\ldots,\alpha_4$ be its roots. Then
\[
\sum_i\alpha_i=-\frac ba,\qquad
\sum_{i<j}\alpha_i\alpha_j=\frac1a,\qquad
\sum_{i<j<k}\alpha_i\alpha_j\alpha_k=-\frac1a,\qquad
\prod_i\alpha_i=\frac1a.
\]
Put $\beta_i=1/\alpha_i$. Then $\sum_i\beta_i=(-1/a)/(1/a)=-1$ and
$\sum_{i<j}\beta_i\beta_j=1$. Thus
\[
\sum_{i=1}^4\beta_i^2=\left(\sum_{i=1}^4\beta_i\right)^2
-2\sum_{i<j}\beta_i\beta_j=(-1)^2-2=-1<0.
\]
One of the $\beta_i$, and therefore one of the $\alpha_i$, must be complex.

\setcounter{theorem}{49}\begin{lemma}[Polynomial agreement]
If $f$ is a polynomial and $f(x)=g(x)$ for $x\in\{1,2,\ldots,k\}$, then
$f(x)-g(x)=(x-1)(x-2)\cdots(x-k)h(x)$.
Moreover, if $f$ has degree $k$, then
$f(x)-g(x)=k'(x-1)(x-2)\cdots(x-k)$, where $k'$ is its leading coefficient.
\end{lemma}
% Author check: no degree/polynomial hypotheses on g are stated.
Let $f$ be a monic polynomial of degree $5$, with
$f(2)=3$, $f(3)=8$, $f(4)=15$, $f(5)=24$, $f(6)=35$. Find $f(1)$.
Observe that $f(x)=x^2-1$ at $x=2,3,4,5,6$.
Thus $f(x)-(x^2-1)=(x-2)(x-3)(x-4)(x-5)(x-6)$, giving $f(1)=-120$.

\setcounter{theorem}{56}\begin{lemma}[Choose useful notation]
Proper notation or a clever choice of variables can turn a complex problem into child's play.
Let me illustrate exactly what I mean.
\end{lemma}
If $x^4+1/x^4$ and $x^5+1/x^5$ are rational for real $x$, prove that $x+1/x$ is rational.
One can work with the powers of $x$ and $1/x$, but notation makes manipulation easy.
Let $T_k=x^k+1/x^k$. Observe
$T_mT_n=T_{m+n}+T_{m-n}$ and $T_{2k}=T_k^2+2$.
The purpose is to show $T_1\in\mathbb Q$.
From $T_4T_2=T_2+T_6$ and $T_8T_2=T_6+T_{10}$,
we get $T_2(T_8-T_4+1)=T_{10}$, so $T_2,T_6$ are rational.
% Source page 26 - left
Finally $T_1T_5=T_6+T_4$, so $T_1$ is rational.
% Author check: T_{2k}=T_k^2+2 is the source sign; denominator nonzero is not discussed.

\setcounter{theorem}{57}\begin{lemma}[Lagrange interpolation]
Let $x_0,x_1,\ldots,x_n$ be distinct real numbers and $y_0,y_1,\ldots,y_n$ arbitrary real numbers.
There exists a unique polynomial $P$ of degree at most $n$ with $P(x_i)=y_i$ for every $i$.
It is given by
\[
P(x)=\sum_{i=0}^n y_i
\frac{(x-x_0)\cdots(x-x_{i-1})(x-x_{i+1})\cdots(x-x_n)}
{(x_i-x_0)\cdots(x_i-x_{i-1})(x_i-x_{i+1})\cdots(x_i-x_n)}.
\]
\end{lemma}
Sometimes knowing such interpolation formulas serves the purpose.

Have a look at these problems:
\begin{enumerate}
\setcounter{enumi}{0}\item Find all integers $a,b,c$ such that
$(2a+b-2)^3+(2b+a-1)^3+27(1-a-b)^3=2014$.
\setcounter{enumi}{1}\item Prove that $\sqrt[3]{2+\sqrt5}+\sqrt[3]{2-\sqrt5}$ is rational.
\end{enumerate}
My usual practice is to state a lemma before its application.
This time a glance at the lemma would give the solution; do not peek down when seeking the solution!
\setcounter{theorem}{62}\begin{lemma}[Three cubes with zero sum]\label{lemma:lb1-threecubes}
If $a+b+c=0$, then $a^3+b^3+c^3=3abc$.
In general,
$a^3+b^3+c^3=(a+b+c)(a^2+b^2+c^2-ab-bc-ac)+3abc$.
\end{lemma}
There are many proofs using complex numbers and other stuff.
Here is my preferred two-line proof without factoring $a^3+b^3+c^3-3abc$.
Recall the identity $(a+b+c)^3=a^3+b^3+c^3+3(a+b)(b+c)(c+a)$.
If $a+b+c=0$, then
$a^3+b^3+c^3=-3(a+b)(b+c)(c+a)=(-3)(-c)(-a)(-b)=3abc$.

\setcounter{theorem}{64}\begin{lemma}[Vi\`ete's relations]
If $x^n+p_1x^{n-1}+p_2x^{n-2}+\cdots+p_n=0$ has roots $\alpha_1,\ldots,\alpha_n$, then
\[
\sum_{i=1}^n\alpha_i=-p_1,\qquad
\sum_{\mathrm{cyc}}\alpha_{i_1}\alpha_{i_2}=p_2,\qquad
\sum_{\mathrm{cyc}}\alpha_{i_1}\cdots\alpha_{i_k}=(-1)^kp_k,\qquad
\prod_{i=1}^n\alpha_i=(-1)^np_n.
\]
\end{lemma}
% Author check: source labels the elementary symmetric sums ``cyc''; retained.
Given such symmetric expressions as a system, form a polynomial: its roots are the quantities of interest.
% Source page 28 - right
\paragraph{Application.}
Solve $x+y+z=6$, $xy+yz+zx=11$, $xyz=6$.
Generate the polynomial with roots $x,y,z$:
$P(t)=t^3-6t^2+11t-6=(t-1)(t-2)(t-3)$.
Its roots are $(1,2,3)$, which is what we wanted.

\setcounter{theorem}{67}\begin{lemma}[Number of polynomial roots]
A polynomial of degree $n$ has at most $n$ roots, and exactly $n$ roots in the complex domain.
\end{lemma}
\paragraph{Applications: USAMO, rephrased.}
\begin{enumerate}
\setcounter{enumi}{0}\item Let $P$ be a polynomial of degree $n\geq1$, with $P(x)=x/(x+1)$ for $x=1,2,\ldots,n+2$.
Prove that no such polynomial exists.
Set $Q(x)=(x+1)P(x)-x$. It has degree $n+1$ but roots $1,2,\ldots,n+2$, a contradiction.
\setcounter{enumi}{1}\item Let $P:\mathbb N_0\to\mathbb R$ be a polynomial of degree $n\geq1$,
with $P(x)=x/(x+1)$ for $x=0,1,\ldots,n$. Find all possible values of $P(2014)$.
\end{enumerate}

\subsection{Identity principle and inequalities}
\setcounter{theorem}{68}\begin{lemma}[Identity principle]\label{lemma:lb1-72}
If $\deg(P)=\deg(Q)<n$ and $P(x)=Q(x)$ for $n$ values, then $P\equiv Q$.
\end{lemma}
Use Lemma 61 in the manuscript to prove the principle.
% The source explicitly refers to its Lemma 61 (subtracting two polynomials).

\paragraph{An ingenious technique to factorize: Euler's enigma.}
Factor $ab(a^2-b^2)+bc(b^2-c^2)+ca(c^2-a^2)$.
Let
\[P(t)=tb(t^2-b^2)+bc(b^2-c^2)+tc(c^2-t^2).\]
Observe that $P(b),P(c),P(-b-c)$ all equal zero. ($P$ is a biquadratic polynomial.)
Thus $P(t)=K(t-b)(t-c)(t+b+c)$.
Now make another crucial observation:
\[
P(t)=t[bt^2-b^3+c^3-ct^2]+bc(b^2-c^2)
=t[t^2(b-c)-(b^3-c^3)]+bc(b^2-c^2).
\]
Hence $b-c\mid P(t)$, so $K$ must be $b-c$; it is the only linear factor we missed!
% Source page 32 - right
Thus $|P(t)|=|(b-c)(t-b)(t-c)(t+b+c)|$ and
\[
P(a)=|(b-c)(a-b)(a-c)(a+b+c)|.
\]
Therefore
\[\left|ab(a^2-b^2)+bc(b^2-c^2)+ca(c^2-a^2)\right|
=\left|(b-c)(a-b)(a-c)(a+b+c)\right|.\]
Q.E.D.!
% Author check: source calls this polynomial biquadratic and its displayed final formulas use absolute-value bars.
The reason for calling it Euler's enigma is the following note, which he scribbled on his paper!
Let us compute the series $1+1/2^2+1/3^2+\cdots$.
The Taylor expansion of the sine function is
\[\sin x=x-\frac{x^3}{3!}+\frac{x^5}{5!}-\cdots.\]
The zeros of $\sin x$ are $0,\pm\pi,\pm2\pi,\ldots$.
So Euler claimed the factorization
\[\sin x=x-\frac{x^3}{3!}+\frac{x^5}{5!}-\cdots
=x\left(1-\frac x\pi\right)\left(1+\frac x\pi\right)
\left(1-\frac x{2\pi}\right)\left(1+\frac x{2\pi}\right)\cdots.\]
Divide by $x$ to obtain
\[1-\frac{x^2}{3!}+\frac{x^4}{5!}-\cdots
=\left(1-\frac{x^2}{\pi^2}\right)\left(1-\frac{x^2}{4\pi^2}\right)
\left(1-\frac{x^2}{9\pi^2}\right)\cdots.\]
Equate the coefficients of $x^2$ on both sides:
\[-\frac16=-\frac1{\pi^2}-\frac1{4\pi^2}-\frac1{9\pi^2}-\frac1{16\pi^2}-\cdots,
\qquad \frac{\pi^2}{6}=1+\frac14+\frac19+\frac1{16}+\frac1{25}+\cdots.\]

\paragraph{Ingenious application of Vi\`ete's relations.}
Refer to Lemma 68 in the manuscript. Solve in reals the system
\[a+b=8,\qquad ab+c+d=23,\qquad ad+bc=28,\qquad cd=12.\]
Try to solve the problem independently; that will prepare you to appreciate the solution.
\[
\begin{aligned}
(X^2+aX+c)(X^2+bX+d)
&=X^4+(a+b)X^3+(ab+c+d)X^2+(ad+bc)X+cd=0,\\
X^4+8X^3+23X^2+28X+12&=0
\quad\Longrightarrow\quad (X+1)(X+2)^2(X+3)=0.
\end{aligned}
\]
The product is also
\[(X^2+4X+3)(X^2+4X+4)=(X^2+3X+2)(X^2+5X+6)=0.\]
Hence
\[(a,b,c,d)\in\{(4,4,3,4),(4,4,4,3),(3,5,2,6),(5,3,6,2)\}.\]
Q.E.D.!
% Source page 35 - right
\Needspace{14\baselineskip}
\subsection{A novel way of solving}
\setcounter{theorem}{76}\begin{lemma}[Real roots of a cubic]\label{lemma:lb1-79}
The condition for a cubic equation with real coefficients to have only real roots is derived analytically.
The equation $x^3+ax^2+bx+c=0$, with $a,b,c$ real, has real roots if
\[(2a^3-9ab+27c)^2\leq4(a^2-3b)^3.\]
\end{lemma}
Let us use this to solve a totally unexpected problem!
\paragraph{Application.}
If $x,y,z>0$ and $xy+yz+zx=2(x+y+z)$, prove that $xyz\leq x+y+z+2$.
Here $x,y,z$ are real.
Create a polynomial with roots $x,y,z$:
\[P(X)=X^3-mX^2+\frac m2X-c,\qquad m=x+y+z.\]
$P$ has all real solutions. Let $c>m+2$; this is our assumption!
The above lemma gives
\[\left(2(-m)^3-9(-m)\frac m2+27(-c)\right)^2
\leq4\left(m^2-\frac{3m}{2}\right)^3.\]
Let me skip some algebraic steps. An equivalent form is
% Source page 36 - left
\[108c^2+c(16m^3-36m^2)+2m^3-m^4\leq0.\]
This is an increasing function of $c$, so $c=m+3$, the least value of $c$, must satisfy the inequality.
Substituting $c=m+3$ gives
\[15m^4+14m^3+648m+972\leq0.\]
But $m=x+y+z>0$, so we have reached a contradiction.
Hence $c\leq m+2$, and therefore $xyz\leq x+y+z+2$. Q.E.D.!
You can also resolve the problem using Lagrange multipliers!
% Author check: the source uses m/2 as the quadratic elementary sum although the hypothesis gives 2m. It also uses m+3 as the least value above m+2 despite real variables; the asserted monotonicity is unqualified.

\paragraph{V. V. Acharya's enigmatic problem: see the unseen.}
Let $P(x)=x^8-4x^7+7x^6+ax^5+bx^4+cx^3+dx^2+ex+f$.
The roots of this polynomial are positive reals. Determine the constants $(a,b,c,d,e,f)$.
Suppose the roots are $y_1,y_2,y_3,\ldots,y_8$. By Vi\`ete's relations,
\[\sum y_i=4,\qquad\sum_{i<j}y_iy_j=7,
\qquad\sum y_i^2=\left(\sum y_i\right)^2-2\sum_{i<j}y_iy_j=2.\]
Note that $\mathrm{AM}=(\sum y_i)/8=1/2$, and $\mathrm{RMS}=\sqrt{(\sum y_i^2)/8}=1/2$. Wow!!
Since AM $=$ RMS, we have $y_1=y_2=\cdots=y_8=1/2$.
Now determine the constants for yourself. Q.I.D.: Quite Intelligently Done.

% Source page 53 - left
\subsection{It was obvious!}
Given $a/(b+c)+b/(c+a)+c/(a+b)=1$, find
$a^2/(b+c)+b^2/(c+a)+c^2/(a+b)$.
Don't frown at the condition: Nesbitt's inequality works if $a,b,c$ are nonnegative, which is not necessary here!
We have
\[\frac a{b+c}+\frac b{c+a}+\frac c{a+b}=\frac{a+b+c}{a+b+c},
\qquad (a+b+c)\sum\frac a{b+c}=a+b+c.\]
Thus $a^2/(b+c)+b^2/(c+a)+c^2/(a+b)+(a+b+c)=a+b+c$.
It follows that $\sum_{\mathrm{cyc}}a^2/(b+c)=0$. Q.E.D.!
% Author check: the source introduces (a+b+c)/(a+b+c) without discussing a+b+c=0.

\paragraph{When did the IMO get so easy! IMO shortlist (1967).}
Solve the cyclic system $x_i^2+x_i-1=x_{i+1}$ for $1\leq i\leq n$, where $x_{n+1}=x_1$.
It is not obvious. Note two things:
\[
\sum_{\mathrm{cyc}}(x_i^2+x_i-1)=\sum_{\mathrm{cyc}}x_i
\Longrightarrow\sum_{\mathrm{cyc}}x_i^2=n,
\]
and $x_i(x_i+1)=x_{i+1}+1$, so
$\prod_{\mathrm{cyc}}x_i(x_i+1)=\prod_{\mathrm{cyc}}(x_{i+1}+1)$,
giving $\prod_{\mathrm{cyc}}x_i=1$. Ha ha! There you go.
Then $\sum_i x_i^2/n\geq\prod_i x_i$, so $\prod_i x_i\leq1$ and
$x_1=x_2=\cdots=x_n=\pm1$. That is it!
% Author check: zero factors and AM--GM powers/sign conditions are not treated; retained.
% BLOG-CONTENT-END
\end{document}
